\documentclass[11pt,a4paper]{article}
\usepackage[a4paper,margin=18mm]{geometry}
\usepackage[T1]{fontenc}
\usepackage{amsmath}
\usepackage{amssymb}
\usepackage{enumitem}

\setlist[enumerate,1]{label=\textbf{Q\arabic*.}, leftmargin=*, itemsep=0.8em, topsep=2pt}
\setlist[enumerate,2]{label=(\alph*), leftmargin=*, itemsep=0.25em, topsep=1pt}
\setlength{\parindent}{0pt}
\setlength{\parskip}{3pt}

\begin{document}

\begin{center}
{\Large\bfseries Answer Key: Ratio and the Shape of A4 Paper}
\end{center}

\begin{enumerate}[label=\textbf{Q\arabic*.}, leftmargin=*, itemsep=0.8em]

\item
\begin{enumerate}[label=(\alph*)]
\item A3: \(1:\frac{420}{297}\approx 1:1.414\); \quad
A4: \(1:\frac{297}{210}\approx 1:1.414\); \quad
A5: \(1:\frac{210}{148}\approx 1:1.419\).
\item The ratios are all approximately \(1:\sqrt2\). The sheets are very nearly similar; the small differences arise from rounding to whole millimetres.
\end{enumerate}

\item
\begin{enumerate}[label=(\alph*)]
\item \(1:1.5\) gives half-sheet ratio \(1:1.333\), not similar; \\
\(1:2\) gives \(1:1\), not similar; \\
\(1:1.2\) gives \(1:1.667\), not similar; \\
\(1:1.4\) gives \(1:1.429\), not similar.
\item \(n=\frac2n\), so \(n^2=2\) and \(n=\sqrt2\approx1.414\). This lies between \(1:1.4\) and \(1:1.5\).
\end{enumerate}

\item
\begin{enumerate}[label=(\alph*)]
\item \(10:15\), \ \(1:1.5\), \ \(0.4:0.6\).
\item In each case the second part divided by the first part is \(1.5\).
\item Multiplying or dividing both parts by the same nonzero number leaves \(\frac{b}{a}\) unchanged.
\item \(y:6=1:\frac6y\); \quad \(\frac y4:1=1:\frac4y\).
\end{enumerate}

\item
\(x=\sqrt2\), so the A-series ratio is
\[
1:\sqrt2\approx1:1.414.
\]

\item
Letter sheet:
\[
1:\frac{11}{8.5}\approx1:1.294.
\]
Half Letter sheet:
\[
1:\frac{8.5}{5.5}\approx1:1.545.
\]
These are not equal, so the half sheet is not similar to the whole sheet.

\item
\begin{enumerate}[label=(\alph*)]
\item Short side \(=\frac1{\sqrt[4]{2}}\text{ m}\approx0.8409\text{ m}\approx841\text{ mm}\). \\
Long side \(=\sqrt[4]{2}\text{ m}\approx1.1892\text{ m}\approx1189\text{ mm}\).
\item
\[
\begin{array}{c|ccccccc}
\text{Size} & A0&A1&A2&A3&A4&A5&A6\\ \hline
\text{Short (mm)} & 841&594&420&297&210&148&105\\
\text{Long (mm)} & 1189&841&594&420&297&210&148
\end{array}
\]
\item Four halvings give \(2^4=16\) A4 sheets per A0, so
\[
\text{area of A4}=\frac1{16}\text{ m}^2=625\text{ cm}^2.
\]
Using rounded dimensions, \(21.0\times29.7=623.7\text{ cm}^2\approx625\text{ cm}^2\).
\end{enumerate}

\item
\begin{enumerate}[label=(\alph*)]
\item \(141\%\) turns A4 into A3.
\item \(71\%\) turns A3 into A4.
\item \(200\%\) turns A4 into A2.
\item About \(283\%\) turns A4 into A1.
\end{enumerate}

\item
The fold line is the diagonal of a unit square, so its length is \(\sqrt2\), equal to the long side of an A-series sheet. Thus the paper is A-series.

\item
\[
1:\sqrt3\approx1:1.732.
\]

\item
\begin{enumerate}[label=(\alph*)]
\item \(1:\sqrt5\approx1:2.236\).
\item \(1:\sqrt n\).
\item This is exact with whole numbers when \(\sqrt n\) is an integer, i.e.\ \(n=1,4,9,16,\ldots\). For example, \(n=4\) gives \(1:2\).
\end{enumerate}

\item
\begin{enumerate}[label=(\alph*)]
\item \(1:\frac1{x-1}\).
\item \(x=\frac1{x-1}\), so \(x^2-x-1=0\).
\item \(x=\frac{1+\sqrt5}{2}\approx1.618\). Hence the ratio is
\[
1:\frac{1+\sqrt5}{2}\approx1:1.618.
\]
\end{enumerate}

\item
\begin{enumerate}[label=(\alph*)]
\item \((\sqrt2-1):1\).
\item \((\sqrt2-1):1=1:(\sqrt2+1)\), since \((\sqrt2-1)(\sqrt2+1)=1\).
\item Cutting two unit squares from a \(1:(1+\sqrt2)\) rectangle leaves a \((\sqrt2-1):1\) rectangle, which is similar to the whole.
\end{enumerate}

\item
\begin{enumerate}[label=(\alph*)]
\item \(1:\sqrt2\).
\item \(1:\sqrt3\).
\item \(1:2\).
\item For \(n\) similar equal pieces, the ratio is \(1:\sqrt n\). For A4 halved into two A5 sheets, \(n=2\), giving \(1:\sqrt2\).
\end{enumerate}

\item
\begin{enumerate}[label=(\alph*)]
\item The half-block ratio is
\[
1:\frac2z:\frac{2y}{z}.
\]
\item Similarity gives \(y=\frac2z\) and \(z=\frac{2y}{z}\), hence \(y^3=2\).
\item \(y=\sqrt[3]{2}\approx1.260\), \quad \(z=\sqrt[3]{4}\approx1.587\). Thus
\[
1:\sqrt[3]{2}:\sqrt[3]{4}\approx1:1.260:1.587.
\]
\item \(a=\sqrt[3]{500}\approx7.94\text{ cm}\), \quad middle edge \(=10\text{ cm}\), \quad longest edge \(=\sqrt[3]{2000}\approx12.60\text{ cm}\).
\end{enumerate}

\item
\begin{enumerate}[label=(\alph*)]
\item The half-box ratio is
\[
1:\frac2c:\frac{2a}{c}:\frac{2b}{c}.
\]
\item Comparing gives \(a=\frac2c\), \(b=\frac{2a}{c}\), \(c=\frac{2b}{c}\). Hence
\[
b=a^2,\qquad c=a^3,\qquad a^4=2.
\]
\item
\[
1:a:b:c\approx1:1.189:1.414:1.682.
\]
Also \(a=2^{1/4}\approx1.189\), matching the long side of A0 in metres (\(1189\) mm), and \(b=\sqrt2\), the A-series ratio.
\end{enumerate}

\item
\begin{enumerate}[label=(\alph*)]
\item In general \(r=n^{1/d}\). Thus:
\[
\begin{array}{c|cccc}
 & \text{2D} & \text{3D} & \text{4D} & d\text{ dimensions}\\ \hline
2\text{ slices} & \sqrt2 & \sqrt[3]{2} & 2^{1/4} & 2^{1/d}\\
3\text{ slices} & \sqrt3 & \sqrt[3]{3} & 3^{1/4} & 3^{1/d}\\
n\text{ slices} & \sqrt n & \sqrt[3]{n} & n^{1/4} & n^{1/d}
\end{array}
\]
\item For \(2\) slices, longest divided by shortest is \(2^{1-1/d}\):
\[
\begin{array}{c|c}
2\text{D} & \sqrt2\approx1.414\\
3\text{D} & \sqrt[3]{4}\approx1.587\\
4\text{D} & 2^{3/4}\approx1.682\\
10\text{D} & 2^{9/10}\approx1.866
\end{array}
\]
\item As \(d\to\infty\), \(2^{1-1/d}\to2\), but it is always less than \(2\).
\end{enumerate}

\item
\begin{enumerate}[label=(\alph*)]
\item
\[
\frac12+\frac14+\frac18+\frac1{16}+\frac1{32}+\frac1{64}
=\frac{63}{64}\text{ m}^2.
\]
\item The infinite sum is
\[
\frac12+\frac14+\frac18+\cdots=1\text{ m}^2.
\]
The leftover area tends to \(0\).
\item
\[
\left(\frac1{\sqrt2}\right)^{10}=2^{-5}=\frac1{32}.
\]
So A10 is about
\[
\frac{841}{32}\text{ mm}\times\frac{1189}{32}\text{ mm}
\approx26\text{ mm}\times37\text{ mm}.
\]
\end{enumerate}

\item
\begin{enumerate}[label=(\alph*)]
\item \(x=\frac1{x-k}\), so
\[
x^2-kx-1=0.
\]
\item For \(k=3\),
\[
x=\frac{3+\sqrt{13}}2\approx3.303,
\]
so the ratio is
\[
1:\frac{3+\sqrt{13}}2\approx1:3.303.
\]
\item In general,
\[
x=\frac{k+\sqrt{k^2+4}}2.
\]
For \(k=1\), this is the golden ratio; for \(k=2\), this is \(1+\sqrt2\), the silver ratio.
\item One valid rule is: remove a unit square, split the remaining rectangle into two equal rectangles, and require each small rectangle to be similar to the whole. This gives \(x=2\), so the ratio is \(1:2\). Indeed, a \(1\times2\) rectangle splits into a \(1\times1\) square and two \(\frac12\times1\) rectangles, each similar to the whole.
\end{enumerate}

\end{enumerate}

\end{document}